\clearpage\pagestyle{otMain}\fontsize{11}{15}\selectfont
\corrigesChapitre{12}{Statistique descriptive}
\begin{blocLivre}{cor-12-1}
\begin{corrige}{1}{Moyenne simple}
\demarche{Additionner les valeurs puis diviser par leur nombre.}
La somme est $4+6+7+8+10=35$ pour 5 valeurs. La moyenne est donc $\res{35/5=7}$. Elle se situe entre le minimum 4 et le maximum 10, ce qui est un premier contrôle de cohérence.
\end{corrige}
\end{blocLivre}
\begin{blocLivre}{cor-12-2}
\begin{corrige}{2}{Médiane et quartiles}
\demarche{Ordonner puis utiliser les rangs adaptés à l’effectif.}
Les valeurs centrales sont 4 et 5 : la médiane vaut $\res{4{,}5}$. Les rangs des quartiles sont 2 et 6 : $\res{Q_1=2}$ et $\res{Q_3=6}$. L’étendue vaut $8-1=\res7$. Les conventions de rang doivent être précisées pour éviter les différences entre certains logiciels.
\end{corrige}
\end{blocLivre}
\begin{blocLivre}{cor-12-3}
\begin{corrige}{3}{Une moyenne pondérée}
\demarche{Pondérer chaque valeur par son nombre d’occurrences.}
L’effectif total est $1+3+4+2=10$. La somme pondérée vaut $2\times1+4\times3+6\times4+8\times2=54$. Ainsi $\res{\bar x=5{,}4}$. Faire la moyenne simple de 2, 4, 6 et 8 aurait donné 5 et aurait oublié leurs effectifs différents.
\label{fig-cor-12-3}
\figureGeo[corrige]{stats-barres}
\end{corrige}
\end{blocLivre}
\begin{blocLivre}{cor-12-4}
\begin{corrige}{4}{Lire un secteur}
\demarche{Rapporter l’angle du secteur à l’angle qui représente le total.}
Le cercle entier vaut $360^\circ$, donc la proportion est $120/360=\res{1/3}$ et l’effectif $\res{20\text{ élèves}}$. Dans le demi-cercle, la référence n’est plus 360 mais 180 degrés : $120/180=\res{2/3}$, donc $\res{40\text{ élèves}}$. L’angle seul ne donne pas le pourcentage ; il faut connaître l’angle qui représente toute la population.
\label{fig-cor-12-4}
\figureGeo[corrige]{stats-secteur}
\end{corrige}
\end{blocLivre}
\begin{blocLivre}{cor-12-5}
\begin{corrige}{5}{Un écart type court}
\demarche{Calculer la moyenne des carrés des écarts avant de prendre la racine.}
La moyenne vaut $(2+2+6+6)/4=4$. Les écarts à la moyenne sont $-2,-2,2,2$ ; leurs carrés valent tous 4. La moyenne de ces carrés vaut 4, et l’écart type vaut $\res{\sigma=2}$. L’écart type n’est pas la moyenne des écarts signés, qui est ici nulle.
\end{corrige}
\end{blocLivre}
\begin{blocLivre}{cor-12-6}
\begin{corrige}{6}{Même moyenne, dispersion différente}
\demarche{Comparer un couple position-dispersion plutôt que la seule moyenne.}
Les deux moyennes valent 4. Pour $A$, tous les écarts sont nuls, donc $\sigma_A=0$. Pour $B$, les carrés des écarts sont 9, 1, 1 et 9 : leur moyenne vaut 5, donc $\res{\sigma_B=\sqrt5}$. La série $B$ est plus dispersée autour de la même moyenne.
\label{fig-cor-12-6}
\figureGeo[corrige]{stats-dispersion}
\end{corrige}
\end{blocLivre}
\begin{blocLivre}{cor-12-7}
\begin{corrige}{7}{Transformer une série}
\demarche{Comparer les écarts avant et après une transformation affine.}
Après ajout de 5, la moyenne vaut $\res{17}$. Chaque écart devient $(x_i+5)-(12+5)=x_i-12$ ; l’écart type reste $\res3$. Pour la seconde transformation appliquée à la série initiale, la moyenne vaut $-2\times12+1=\res{-23}$. Les écarts sont multipliés par $-2$, donc leurs carrés par 4 et l’écart type par 2 : il vaut $\res6$. Le signe négatif inverse l’ordre des valeurs sans rendre la dispersion négative.
\end{corrige}
\end{blocLivre}
\begin{blocLivre}{cor-12-8}
\begin{corrige}{8}{Réunir deux groupes}
\demarche{Retrouver les sommes de groupe avant de réunir les effectifs.}
Les sommes des deux groupes valent $12\times8=96$ et $18\times13=234$. La somme totale vaut 330 pour 30 mesures : $\res{\bar x=11}$. La moyenne simple 10,5 serait fausse, car le groupe dont la moyenne est 13 est plus nombreux.
\end{corrige}
\end{blocLivre}
\begin{blocLivre}{cor-12-9}
\begin{corrige}{9}{Moyenne par classes}
\demarche{Utiliser les centres et annoncer la nature de l’estimation.}
Les centres sont 5, 15, 25 et 35 minutes. La moyenne estimée est $(4\times5+8\times15+6\times25+2\times35)/20=\res{18\text{ min}}$. Cette estimation est exacte dans un modèle uniforme dans chaque classe. Les seuls effectifs ne donnent pas la position exacte des valeurs à l’intérieur des intervalles.
\label{fig-cor-12-9}
\figureGeo[corrige]{stats-histogramme}
\end{corrige}
\end{blocLivre}
\begin{blocLivre}{cor-12-10}
\begin{corrige}{10}{Classe médiane et interpolation}
\demarche{Localiser le seuil de 50 \% puis interpoler uniquement sous l’hypothèse annoncée.}
L’effectif total est 20. Les cumuls sont 4, 12, 18 et 20 ; le passage de 10 observations se fait dans $[10;20[$. Pour atteindre 10 depuis le cumul 4, il faut 6 des 8 observations de la classe, soit $6/8$ de sa largeur. L’estimation est donc $10+(6/8)\times10=\res{17{,}5\text{ min}}$.
\label{fig-cor-12-10}
\figureGeo[corrige]{stats-cumuls}
\end{corrige}
\end{blocLivre}
\begin{blocLivre}{cor-12-11}
\begin{corrige}{11}{Un point extrême}
\demarche{Réordonner et mettre à jour à la fois les effectifs et les rangs.}
Avant remplacement, moyenne et médiane valent 4. Après remplacement, la somme vaut 114 et la moyenne $\res{22{,}8}$, tandis que la médiane reste $\res4$. Après suppression de 100, la série $2,3,4,5$ a pour moyenne $14/4=\res{3{,}5}$ et pour médiane $(3+4)/2=\res{3{,}5}$. Une modification d’effectif peut déplacer aussi les rangs centraux : la médiane n’est pas toujours inchangée.
\end{corrige}
\end{blocLivre}
\begin{blocLivre}{cor-12-12}
\begin{corrige}{12}{Deux boîtes à moustaches}
\demarche{Lire uniquement les indicateurs contenus dans une boîte.}
Les deux médianes valent 6. L’écart interquartile de $A$ vaut $8-4=4$, celui de $B$ vaut $7-5=2$. Les étendues valent respectivement 12 et 8. Selon ces deux mesures, $B$ est moins dispersée. Les boîtes ne suffisent pas à connaître les moyennes ni les écarts types exacts.
\label{fig-cor-12-12}
\figureGeo[corrige]{stats-boites}
\end{corrige}
\end{blocLivre}
\begin{blocLivre}{cor-12-13}
\begin{corrige}{13}{L’effet d’un ajout sur la moyenne}
\demarche{Mettre à jour à la fois la somme et l’effectif.}
La somme initiale vaut 140. Après ajout, la somme vaut 165 et l’effectif 11. La nouvelle moyenne est $\res{15}$. Comme 25 est au-dessus de l’ancienne moyenne, l’ajout tire la moyenne vers le haut. Il ne faut pas conserver le dénominateur 10 après avoir ajouté une observation.
\end{corrige}
\end{blocLivre}
\begin{blocLivre}{cor-12-14}
\begin{corrige}{14}{Moyennes de classes connues}
\demarche{Distinguer l’information exacte de l’hypothèse de répartition uniforme.}
La moyenne exacte est $(8\times7+12\times16)/20=\res{12{,}4\text{ min}}$. L’estimation par centres donne $(8\times5+12\times15)/20=\res{11\text{ min}}$. Les moyennes réelles de classes diffèrent des centres : l’hypothèse qui rendrait l’estimation exacte n’est pas satisfaite. Lorsqu’elles sont connues, il faut utiliser les moyennes réelles.
\end{corrige}
\end{blocLivre}
\begin{blocLivre}{cor-12-15}
\begin{corrige}{15}{Une série à construire}
\demarche{Construire des valeurs symétriques autour du centre choisi.}
On peut choisir $A=(4,4,6,6)$ et $B=(0,0,10,10)$. Chaque somme vaut 20 et chaque paire centrale a pour moyenne 5. Pour $A$, les écarts au carré valent tous 1 : $\sigma_A=\res1$. Pour $B$, ils valent tous 25 : $\sigma_B=\res5$. Même moyenne et même médiane n’imposent donc pas la même dispersion.
\end{corrige}
\end{blocLivre}
